\section{Certificate for the scalar enclosures}
\label{rcut:sec:scalar-certificate}

The scalar argument needs certified moments of the field law
\(N(\beta^2,\beta^2)\), with \(\beta\in[7/20,1/2]\). Retain the notation
of Section~\ref{rcut:sec:scalar}. We now prove
Lemma~\ref{rcut:lem:scalar-data} and verify the arithmetic in
Table~\ref{rcut:tab:scalar-intervals}.

There are three distinct error estimates. First, a complex-strip bound
controls replacement of each integral by a finite quadrature sum.
Second, directed integer arithmetic encloses those sums, which use
41 nodes at each of five exact temperatures. Third, normalization,
standard deviations, and the two mean parameters in \(F\) must be
transferred back to the continuous field law. The estimates below keep
these steps separate. The interpolation between temperatures remains the
analytic argument of Section~\ref{rcut:sec:scalar}.

\begin{lemma}[Whole-line trapezoidal estimate]
\label{rcut:lem:scalar-trapezoid}
Let \(t>0\).
Suppose \(f\) is continuous on \(\{z:|\Im z|\leq t\}\), holomorphic
in its interior, and satisfies a bound
\[
 |f(x+iy)|\leq C(1+|x|)^p e^{-c x^2},\qquad |y|\leq t,
\]
for some \(C,c>0\) and \(p\geq0\). If both horizontal-edge integrals
of \(|f|\) are at most \(M\), then, for any mesh \(\Delta>0\),
\begin{equation}
 \left|\Delta\sum_{k\in\mathbb Z}f(k\Delta)
                  -\int_{\mathbb R}f(x)\,dx\right|
 \leq\frac{2M}{e^{2\pi t/\Delta}-1}.
 \label{rcut:eq:scalar-trapezoid-error}
\end{equation}
\end{lemma}

\begin{proof}
We compare the integral and the infinite quadrature sum through the
Fourier coefficients of the periodization. The assumed decay makes
\(\sum_k f(x+k\Delta)\) uniformly convergent on a period. Its
coefficient at integer \(j\) is
\(\Delta^{-1}\int_{\mathbb R}f(x)e^{-2\pi ijx/\Delta}\,dx\).
To estimate a nonconstant coefficient, shift the contour to height
\(-t\) when \(j>0\), and to height \(t\) when \(j<0\). The integrals
on the vertical sides vanish by the uniform Gaussian decay. One may
first shift to an interior line and then pass to the boundary by
dominated convergence, so the stated boundary regularity suffices.
The resulting bound for the absolute value of the coefficient is
\((M/\Delta)e^{-2\pi t|j|/\Delta}\).
These bounds are summable. To identify the Fourier series with the
continuous periodization without a further regularity assumption, take
arithmetic means of the partial Fourier sums. Relative to normalized
measure on a period, their kernels are
\(N^{-1}|\sum_{j=0}^{N-1}e^{2\pi ijx/\Delta}|^2\).
The kernels are nonnegative, have integral one, and concentrate at the
origin. Their convolutions therefore converge uniformly to the
continuous periodization. Absolute convergence identifies this limit
with the Fourier series itself. At zero, the constant coefficient is
the integral divided by \(\Delta\); summing the bounds for the remaining
coefficients proves Equation~\eqref{rcut:eq:scalar-trapezoid-error}.
\end{proof}

\paragraph{A common quadrature error for all moments.}
To apply the trapezoidal estimate uniformly in the temperature, change
variables to \(x=h/\beta\), whose density is
\(\varphi(x-\beta)=(2\pi)^{-1/2}e^{-(x-\beta)^2/2}\).
On \(|\Im x|\leq1.5\), we have \(|\Im(\beta x)|\leq.75<\pi/4\).
The identity
\[
 |\tanh(u+it)|^2
  =\frac{\sinh^2u+\sin^2t}{\sinh^2u+\cos^2t}
\]
therefore implies
\begin{equation}
 |m|\leq1,\qquad |v|\leq2,\qquad |\Lambda{}|\geq.14,
 \qquad |1.75-m|\geq.75,\qquad |2-m|\geq1.
 \label{rcut:eq:scalar-strip-bounds}
\end{equation}
For \(d(1)\), first cancel the denominator and use the expression
\(v(1+m)\). The displayed bounds then keep every remaining denominator
in the moment integrands away from zero. The integrands are holomorphic
on a neighborhood of the closed strip and, after multiplication by the
Gaussian density, satisfy the uniform decay in
Lemma~\ref{rcut:lem:scalar-trapezoid}.

The function \(F\) involves two expectations, so its quadrature uses
approximations to those expectations. We need the same analytic bound
for every possible pair of means that will arise. Let \(a',G'\) be
arbitrary constants in \([0,1]\), put \(l'=.64-.19a'\), and let
\(F_{a',G'}\) be the expression in Equation~\eqref{rcut:eq:scalar-F}
with these constants. Since \(l'\geq.45\),
\(|.5-\beta^2a'/l'|\leq.5\). Equation~\eqref{rcut:eq:scalar-strip-bounds}
then gives, with \(h=\beta x\),
\begin{equation}
 |F_{a',G'}(\beta x)|\leq1.5|x|+4.12.
 \label{rcut:eq:scalar-strip-F}
\end{equation}
On either horizontal edge, this is at most \(1.5|\Re x|+6.37\).
The Gaussian modulus on that edge is
\(e^{1.125}\varphi(\Re x-\beta)\). Hence its product with
\(|F_{a',G'}|^2\) has integral at most
\[
 e^{1.125}\bigl(1.5\sqrt{1.25}+6.37\bigr)^2<260.
\]
For a simple upper estimate in the last inequality, use
\(e^{1.125}<4\) and \(\sqrt{1.25}<1.12\).
We also need the raw moments of \(v,g,mv,\rho\), the squares of
\(v,g,mv\), and the five scaled integrands in
Equation~\eqref{rcut:eq:scalar-five-enclosures}. Their edge integrals
satisfy smaller bounds. For example, \(|g|\leq8/3\),
\(|(1+m)v/2-1.5g^2|\leq38/3\), and
\(|(1+v)^2/(16\Lambda{})|\leq9/(16\cdot.14)\).
Consequently \(M=1000\) works simultaneously for every first moment,
raw second moment, and total-mass integral needed below.

The whole-line estimate leaves a second task: bounding the nodes
discarded from the infinite sum. Choose mesh \(\Delta=.4\) and retain
\(|x|\leq8\). At the omitted nodes \(|x|\geq8.4\), every relevant
real moment integrand is bounded by \((1.5|x|+4.12)^2\).
Since \(|\beta|\leq.5\), the sum of the two tails is bounded by twice
the series with terms
\[
 \frac{.4}{\sqrt{2\pi}}(16.72+.6j)^2
                e^{-(7.9+.4j)^2/2},\qquad j=0,1,\ldots.
\]
The ratio of successive terms is at most
\[
 q_*:=e^{-3.24}\left(\frac{17.32}{16.72}\right)^2<.042026.
\]
The two omitted tails consequently have total at most
\begin{equation}
 \frac{.8(16.72)^2 e^{-31.205}}{\sqrt{2\pi}(1-q_*)}
 <2.612\cdot10^{-12}.
 \label{rcut:eq:scalar-tail-error}
\end{equation}
Lemma~\ref{rcut:lem:scalar-trapezoid} bounds the whole-line error by
\(2000/(e^{7.5\pi}-1)<1.171\cdot10^{-7}\).
Together with Equation~\eqref{rcut:eq:scalar-tail-error}, this proves the
following conservative common bound for each unnormalized raw moment:
\begin{equation}
 E_{\mathrm{quad}}:=2.3\cdot10^{-7}.
 \label{rcut:eq:scalar-raw-error}
\end{equation}
This error includes both the infinite-sum error and the omitted tails;
it precedes normalization of the quadrature weights. Each decimal
comparison used to obtain it is checkable by the directed recipe below.
The coarse enclosure
\(3.14159<\pi<3.14160\) suffices; it follows, for example, from
\(\pi=16\arctan(1/5)-4\arctan(1/239)\) and ten terms of the
alternating series for each arctangent.

\paragraph{The finite quantities to be enclosed.}
We next specify the sums exactly, so that their arithmetic verification
is a finite calculation with rational inputs. For \(-20\leq j\leq20\), put
\begin{equation}
 \begin{split}
 x_j&=2j/5,\qquad h_j=\beta x_j,
 \qquad c_j=e^{-(x_j-\beta)^2/2},\\
 y_j&=\frac{e^{2h_j}-1}{e^{2h_j}+1},\qquad
 V_j=1-y_j^2,\qquad P_j=\frac{V_j}{7/4-y_j},\\
 N(f)&=\frac{\sum_{j=-20}^{20}c_j f_j}{\sum_{j=-20}^{20}c_j},
 \qquad T(f)=N(f^2)-N(f)^2.
 \end{split}
 \label{rcut:eq:scalar-finite-sums}
\end{equation}
Define \(a_N=N(V)\), \(G_N=N(P)\), and
\[
 f_j=(V_j+a_N)h_j+2\beta^2y_jV_j+
       \left(.5-\frac{\beta^2a_N}{.64-.19a_N}\right)V_j
       +.12y_j+3\beta^2G_NP_j.
\]
In the order of the columns of Table~\ref{rcut:tab:scalar-moments},
the eight quantities computed from these weights are
\begin{equation}
 N(V),\ N(P),\ \sqrt{T(V)},\ \sqrt{T(P)},\ N(yV),\
 N\!\left(\frac{V(1+y)}2-1.5P^2\right),\ \sqrt{T(f)},\
 \sqrt{T(yV)}.
 \label{rcut:eq:scalar-eight-sums}
\end{equation}
The directed recipe below encloses each of these quantities within
\(10^{-6}\) of the corresponding entry in
Table~\ref{rcut:tab:scalar-finite-values}. The five temperatures in that
table are exact. These enclosures in particular place every finite
quantity within \(.00065\) of the center in
Table~\ref{rcut:tab:scalar-moments}, and place every finite standard
deviation above \(.08\). The latter lower bound will allow us to
transfer variance errors to standard-deviation errors.

\begin{table}[hbp]
 \centering\small
 \caption{Finite-sum centers with radius \(10^{-6}\); the values of
 \(\beta\) are exact. Columns follow
 Equation~\eqref{rcut:eq:scalar-eight-sums}.}
 \label{rcut:tab:scalar-finite-values}
 \setlength{\tabcolsep}{3pt}
 \begin{tabular}{crrrrrrrr}
 \toprule
 \(\beta\)&\(N(V)\)&\(N(P)\)&\(\sqrt{T(V)}\)&\(\sqrt{T(P)}\)&
 \(N(yV)\)&\(N(\rho)\)&\(\sqrt{T(f)}\)&\(\sqrt{T(yV)}\)\\
 \midrule
 .350&.890140&.553815&.128292&.084412&.081332&.014982&.659034&.237553\\
 .400&.860648&.545767&.155504&.093100&.095751&.018407&.729518&.247528\\
 .445&.832336&.536920&.179354&.101479&.107385&.021988&.786378&.252766\\
 .475&.812698&.530184&.194688&.107439&.114316&.024549&.820670&.254625\\
 .500&.795946&.524071&.207033&.112633&.119553&.026744&.846977&.255330\\
 \bottomrule
 \end{tabular}
\end{table}

The five additional endpoint enclosures use the same weights at
\(\beta=1/2\). Their finite expressions are
\begin{equation}
 \begin{gathered}
 \frac1{16}N\!\left(\frac{(1+V)^2}{.45+.25y^2-.06y}\right),\qquad
 \frac1{16}N\!\left(\frac{V^2}{.45+.25y^2-.06y}\right),\\
 \frac14N\!\left(\frac{V^2}{2-y}\right),\qquad
 \frac14N\!\left(\frac{V^2}{(2-y)^2}\right),\qquad
 \frac14N\bigl(V(1+y)\bigr).
 \end{gathered}
 \label{rcut:eq:scalar-five-sums}
\end{equation}
They lie, respectively, within \(10^{-6}\) of
\begin{equation}
 .427718,\quad .089859,\quad .092577,\quad .052281,\quad .228875.
 \label{rcut:eq:scalar-five-finite-values}
\end{equation}

\paragraph{Directed integer arithmetic.}
We now give the finite recipe that establishes the enclosures just
listed. It requires no accuracy assumption for a transcendental
library: exponentials and square roots are enclosed using integer
arithmetic and explicit remainders.
Fix \(S=10^{50}\), and represent an interval by integer endpoints
\([u,v]\), meaning \([u/S,v/S]\). Represent a rational input by its
lower and upper nearest multiples of \(S^{-1}\). Addition and
subtraction act on endpoints in the usual way. For multiplication,
take the minimum and maximum of the four integer endpoint products,
then divide by \(S\), rounding the minimum down and the maximum up.
For an interval not containing zero, its reciprocal is represented by
\[
 \left[\left\lfloor\frac{S^2}{v}\right\rfloor,
        \left\lceil\frac{S^2}{u}\right\rceil\right].
\]
The reciprocal rule applies to positive and negative intervals alike;
division is multiplication by this reciprocal. For a nonnegative
interval, square roots are represented
by \([\lfloor\sqrt{uS}\rfloor,\lceil\sqrt{vS}\rceil]\), computed by
integer square root and an exact square comparison. Powers use
repeated interval multiplication. Every operation preserves inclusion.

All exponential arguments \(d\) needed in
Equations~\eqref{rcut:eq:scalar-finite-sums}--\eqref{rcut:eq:scalar-five-sums}
satisfy \(|d|\leq36.125\). For \(z=d/32\), compute
\[
 p_{20}(z)=\sum_{i=0}^{20}\frac{z^i}{i!}
\]
with the preceding interval operations, for example by the recursion
\(t_0=1\), \(t_i=t_{i-1}z/i\), followed by summation. Enlarge its
interval on both sides by an outward enclosure of
\[
 \frac{4(1.13)^{21}}{21!},
\]
and raise this interval to the 32nd power. This encloses \(e^d\):
\(|z|<1.13\), and Taylor's remainder is at most the displayed
quantity because \(e^{1.13}<4\). In particular the relative error
of \(p_{20}(d/32)^{32}\) is bounded by
\begin{equation}
 \left(1+\frac{16(1.13)^{21}}{21!}\right)^{32}-1<10^{-12}.
 \label{rcut:eq:scalar-exponential-error}
\end{equation}
This relative-error estimate concerns the polynomial approximation
\(p_{20}(d/32)^{32}\). The interval enclosure follows directly from
inserting the absolute remainder before taking the 32nd power, together
with outward rounding at every operation. With
the rational enclosure for \(\pi\) given above, the same recipe checks
the constants in Equations~\eqref{rcut:eq:scalar-tail-error} and
\eqref{rcut:eq:scalar-raw-error} and the exponential in
Equation~\eqref{rcut:eq:scalar-interval-cross}.

Apply these operations to Equations~\eqref{rcut:eq:scalar-finite-sums}--
\eqref{rcut:eq:scalar-eight-sums} to obtain
Table~\ref{rcut:tab:scalar-finite-values}; applying them to
Equation~\eqref{rcut:eq:scalar-five-sums} gives
Equation~\eqref{rcut:eq:scalar-five-finite-values}. Applying the same
operations to the rational endpoints in
Equations~\eqref{rcut:eq:scalar-interval-slacks} and
\eqref{rcut:eq:scalar-interval-cross} gives
Table~\ref{rcut:tab:scalar-intervals}. The loop bounds and remainder
estimates therefore reduce the finite verification to integer arithmetic
alone.
The supplementary files \nolinkurl{numerics-interval.py} and
\nolinkurl{numerics-cross-errors.py} implement these operations and the
error transfers below; their use is not needed to specify the
certificate.

\paragraph{From raw sums to normalized moments.}
The finite arithmetic is now specified. It remains to show that its
enclosures imply the continuous-law enclosures used in the scalar proof.
Let \(Q\) be the unnormalized quadrature with weights
\(.4\varphi(x_j-\beta)\), so \(N(f)=Q(f)/Q(1)\).
The common multiplicative factor omitted in
Equation~\eqref{rcut:eq:scalar-finite-sums} cancels in this ratio.
For every raw moment under consideration,
\(|Q(f)-\langle f\rangle|\leq E_{\mathrm{quad}}\) and
\(|Q(1)-1|\leq E_{\mathrm{quad}}\), with \(E_{\mathrm{quad}}\) from
Equation~\eqref{rcut:eq:scalar-raw-error}.
If the true \(L^2\) norm of \(f\) is at most 4, the normalized first-
and second-moment errors are consequently bounded by
\[
 E_1=\frac{5E_{\mathrm{quad}}}{1-E_{\mathrm{quad}}},\qquad
 E_2=\frac{17E_{\mathrm{quad}}}{1-E_{\mathrm{quad}}}.
\]
Their variance error is at most
\begin{equation}
 E_{\mathrm{var}}=E_2+E_1(8+E_1)
 <1.312\cdot10^{-5}.
 \label{rcut:eq:scalar-variance-error}
\end{equation}
For each of the computed standard deviations,
\(\sqrt{T(f)}>.08\) by Table~\ref{rcut:tab:scalar-finite-values}.
The identity for a difference of square roots therefore bounds its
fixed-function error by
\begin{equation}
 E_{\mathrm{sd}}=E_{\mathrm{var}}/.08
 <1.64\cdot10^{-4}.
 \label{rcut:eq:scalar-sd-error}
\end{equation}
The denominator here comes from the finite computation itself. In
particular, this argument does not assume the continuous-law enclosures
in Table~\ref{rcut:tab:scalar-moments}, which are still to be proved.

We verify the \(L^2\) hypothesis before applying these transfer bounds.
The functions \(v,g,mv\) have absolute value at most one on the real
line. For the fixed function \(F_{a',G'}\), the required bound follows
from
\[
 |F_{a',G'}(h)|\leq2|h|+1.87,
 \qquad
 \|F_{a',G'}\|_2
 \leq2\sqrt{\beta^2+\beta^4}+1.87<3<4.
\]
The real function \(g\) belongs to \([0,1]\), since
\(1-m^2\leq1.75-m\). Positivity of the quadrature weights gives
\(a_N,G_N\in[0,1]\). It follows that all preceding strip and moment
estimates apply to the fixed function
\(F_N:=F_{a_N,G_N}\), whose values at the quadrature nodes are
the \(f_j\) used in Equation~\eqref{rcut:eq:scalar-eight-sums}.

\paragraph{Replacing the two means in \(F\).}
There are two different functions in the remaining comparison. The
finite calculation uses \(F_N\), whereas the desired continuous
standard deviation is that of \(F\), with the true means \(a,G\).
First apply Equation~\eqref{rcut:eq:scalar-sd-error} to the fixed
function \(F_N\). After this step, both standard deviations to be
compared are under the same continuous field law. We may then vary
the two parameters from \(a_N,G_N\) to \(a,G\).
Because \(v,g\) have absolute values at most one,
\begin{equation}
 |a-a_N|,\ |G-G_N|\leq e_a:=\frac{2E_{\mathrm{quad}}}{1-E_{\mathrm{quad}}}
 <4.61\cdot10^{-7}.
 \label{rcut:eq:scalar-mean-error}
\end{equation}
Throughout the parameter rectangle \([0,1]^2\), the derivatives of
\(F_{a',G'}\) are
\[
 \frac{\partial F_{a',G'}}{\partial a'}
   =h-\frac{.64\beta^2}{(.64-.19a')^2}v,
 \qquad
 \frac{\partial F_{a',G'}}{\partial G'}=3\beta^2g.
\]
Their centered \(L^2\) norms are at most
\(\tfrac12+\tfrac14(.64)/(.45)^2\) and \(3/4\), respectively.
The triangle inequality for the centered seminorm consequently gives
\begin{equation}
 \begin{split}
 |s_F-s_{F_N}|&\leq E_{\mathrm{par}}
 :=\left(\frac12+\frac{.64}{4(.45)^2}+\frac34\right)e_a\\
 &<10^{-6}.
 \end{split}
 \label{rcut:eq:scalar-parameter-error}
\end{equation}
The parameter \(l\) needs no separate perturbation allowance: the
derivative above already includes its dependence \(l=.64-.19a\).
Combining Equations~\eqref{rcut:eq:scalar-sd-error} and
\eqref{rcut:eq:scalar-parameter-error}, the total error between the
computed standard deviation of \(f_j\) and the desired \(s_F\) is
\begin{equation}
 E_{\mathrm{sd}}+E_{\mathrm{par}}
 <.000165.
 \label{rcut:eq:scalar-total-error}
\end{equation}

\begin{proof}[Proof of Lemma~\ref{rcut:lem:scalar-data}]
We now combine the finite enclosures with the transfer estimates.
Every quantity in Table~\ref{rcut:tab:scalar-finite-values} is within
\(.00065\) of its center in Table~\ref{rcut:tab:scalar-moments}.
Among the standard deviations, the largest transfer error is
Equation~\eqref{rcut:eq:scalar-total-error}; the three standard deviations
not involving \(F\) require only
Equation~\eqref{rcut:eq:scalar-sd-error}. For the first moments
\(a,G,\langle mv\rangle,R\), the bounded-integrand error is at most
\(e_a\), since Equation~\eqref{rcut:eq:scalar-residual-diagonal} bounds
the integrand of \(R\) by one as well. Thus
\[
 .00065+.000165=.000815<.001
\]
establishes every claimed interval in
Table~\ref{rcut:tab:scalar-moments}.
To finish the five additional endpoint bounds, observe that every
scaled real integrand in
Equation~\eqref{rcut:eq:scalar-five-sums} is bounded by one, so its
normalization and quadrature error is at most \(e_a\).
Adding this to the finite enclosures in
Equation~\eqref{rcut:eq:scalar-five-finite-values} proves the five strict
intervals in Equation~\eqref{rcut:eq:scalar-five-enclosures}.
All numerical comparisons used in this proof follow from the
directed operations specified above.
\end{proof}
